\documentclass[10pt]{article}
\usepackage[margin=20mm]{geometry}
\usepackage{amsmath,amssymb,mathtools,fontspec,unicode-math,xurl,hyperref}
\IfFontExistsTF{DejaVu Serif}{\setmainfont{DejaVu Serif}}{\setmainfont{Latin Modern Roman}}
\setmathfont{latinmodern-math.otf}
\setlength{\parindent}{0pt}\setlength{\parskip}{6pt}\sloppy
\title{Initial Scientific Review of
Uniqueness, Regularity, and Characteristic Flow
for a Non Strictly Convex Singular Variational Problem}
\author{Bounded scientific review: M25-0429}
\date{}
\begin{document}
\maketitle
Abstract: The uniqueness argument uses the common stress, characteristic superlevel sets and boundary traces under substantive geometric hypotheses. Two corrections are established: an oblique slicing formula omits its Jacobian, and the continuity proof uses convexity beyond what assumption (H) supplies. Explicit calibrated examples verify both printed errors. The common Jacobian cancels in the needed inequality, and local convex sweeps and a harmonic continuation argument repair the continuity steps under the original hypotheses. The slicing defect is Minor. The convexity misuse is Major because it invalidates two essential steps in the original proof of the main displacement-continuity result. The supplied replacement repairs that derivation; no counterexample to the stated uniqueness or continuity conclusions is established. Scientific adjudication remains pending.\par
\section*{Identity and scope}
The reviewed source is J.-F. Babadjian and G. A. Francfort, Uniqueness, regularity, and characteristic flow for a non strictly convex singular variational problem [source], assigned DOI https://doi.org/10.1002/cpa.70015. The acquired 44-page file is arXiv:2312.01849v2: its margin dates the version 24 October 2024 and its title page gives 25 October 2024. Its SHA256 is `05f596d9feb2688249f0340eaa472524ff856d46d4e4436c4fc60ae73d8eae16`. Equivalence with a later journal version has not been established. All page numbers below refer to this acquired PDF.\par
No usable substantive predecessor report was recovered in the bounded preserved-reference search. This is an initial review. The primary executor and a native helper each read the complete source, including the appendix and references. The primary executor read the complete analytical checks and the supplemental continuity argument, rechecked their decisive steps, and inspected the indicated source images. The helper is an internal auxiliary role, not independent acceptance. No author code was executed. Optional code and supplements were unavailable in the recorded bounded search; that is not a mathematical defect. Imported results from the paper’s reference [5] were examined as stated inputs, not independently reacquired or reproved.\par
\section*{Framework and positive checks}
The energy density is quadratic for \(|\xi|\leq1\) and equals \(|\xi|-1/2\) for \(|\xi|\geq1\). Its infimal-convolution representation separates the elastic gradient and the plastic measure. Strict convexity of the dual quadratic problem makes the stress \(\sigma\) unique; it does not by itself make the primal displacement unique. The constraints are \[\operatorname{div}\sigma=0,\qquad |\sigma|\leq1,\qquad
 Du=\sigma\,\mathcal L^2+p,\qquad p=\sigma\,|p|.\] In the elastic region the flow rule gives \(p=0\), so \(\sigma=\nabla u\) and \(u\) is harmonic. Under (H), the saturation set \(\Omega_1\) is convex; this is distinct from convexity of \(\Omega\). The standing nonempty-interior assumption yields the convex open plastic region \(\Omega_{\mathrm{pl}}=\operatorname{int}\Omega_1\). The actual source contains the closure in \(\Omega_{\mathrm{el}}=\Omega\setminus\overline{\Omega_1}\); an OCR omission here is not a finding.\par
The characteristic structure in \(\Omega_{\mathrm{pl}}\) uses the locally Lipschitz unit field \(\sigma\) and its constancy on straight chords in direction \(\sigma^\perp\). The characteristic boundary has at most two connected components, each a segment or a singleton. The stress-continuity result allows at most two exceptional interface points. The displacement result concerns \(\Omega_{\mathrm{el}}\cup\Sigma\cup\omega\), with \(\omega\) the region swept from the interface; it is not a general assertion of continuity throughout every plastic component.\par
Theorem 5.1 additionally assumes a simply connected \(C^{1,1}\) domain, pure Dirichlet data, nonempty saturation and (H). In its elastic part, two minimizers have the same gradient; boundary trace attainment on a set of positive length fixes the remaining constant. The boundary flux/perimeter argument in Lemma 5.2 excludes the degenerate alternative. In the plastic region, the characteristic and slicing arguments make superlevel sets half-plane intersections. A difference of two such superlevel sets forces a positive-length boundary level set. This can occur for only countably many levels of the fixed \(L^1\) boundary datum. Integration in the level variable then proves equality almost everywhere. This mechanism and its trace hypotheses were checked; it does not follow from strict convexity of the primal energy.\par
\section*{Problems requiring correction}
Problem 1 (P01: Minor—the oblique slicing identity omits a Jacobian).  Equation (3.16), p. 15, is written without a geometric factor, and the following paragraph asserts the same formula for nonorthogonal slices. Let \(\eta\) be a unit tangent to the slice base line, let \(\xi_0\) be its unit normal, and let \(\xi\) be the unit slicing direction. For \[F(s,t)=x_0+s\eta+t\xi,\qquad
 J=|\det(\eta,\xi)|=|\xi_0\cdot\xi|>0,\] Lebesgue measure pulls back to \(J\,ds\,dt\). Consequently the tensor representation of \(Du\cdot\xi\) requires the same factor \(J\): \[F^*(Du\cdot\xi)=J\,ds\otimes D_t(u\circ F),\] where this notation means equality of the associated measures in the slice coordinates.\par
This is a genuine error on an admissible calibrated solution. Take \(\Omega=(-3,3)^2\), \(u=w=x_1\), \(\sigma=e_1\), \(p=0\), \(\eta=e_2\) and \(\xi=(3/5,4/5)\). Then \(F(s,t)=(3t/5,s+4t/5)\) and \(J=3/5\). On the parallelogram \(P=F((0,1)^2)\), \[(Du\cdot\xi)(P)=\frac35\,|P|=\frac9{25},
 \qquad
 \int_0^1 D_t(u\circ F(s,\cdot))((0,1))\,ds=\frac35.\] The stress satisfies the equilibrium, yield and flow conditions, and its saturation region is the whole convex domain. Thus the example also respects the nonempty plastic setting.\par
The correction is local. The disintegration of \((\sigma\cdot\xi)\mathcal L^2\) carries the identical positive factor \(J\). It cancels when the proof compares the two measures on the selected positive-stress slices. The needed one-dimensional monotonicity and the half-plane description therefore survive. Equation (3.18)’s extra dot with \(\xi\) after a scalar slice derivative is likewise a local notation correction. Neither error supplies a counterexample to the superlevel or uniqueness conclusion.\par
Problem 2 (P02: Major—ambient convexity is used beyond (H)).  In Lemma 4.12, p. 28, the set \[A_\delta=\Omega_{\mathrm{pl}}\cap\{x:\operatorname{dist}(x,\partial\Omega)>\delta\}\] is called convex and is used as the domain in Merlet’s lemma. Assumption (H) makes \(\Omega_{\mathrm{pl}}\) convex, but does not make this ambient distance truncation convex. The later invocation of Proposition 3.25 on p. 29 has the same hypothesis problem: that proposition explicitly assumes convexity of \(\Omega\), whereas Theorem 4.10 assumes (H).\par
A calibrated counterexample to the first geometric claim. Start with the upper half-disc \(|z|<3\), \(\operatorname{Im}z>0\). Write \(r=|z|\), \(\theta=\arg z\), \(R=|z+1|\) and \(\phi=\arg(z+1)\). The following restrictions of the classical crack field are explicit: \[\begin{array}{lll}
r>1:&u=2\sqrt r\,\sin(\theta/2),&
 \sigma=r^{-1/2}(-\sin(\theta/2),\cos(\theta/2)),\\[2pt]
r<1:&u=2\sin\phi,&
 \sigma=(-\sin\phi,\cos\phi).
\end{array}\] In the plastic half-disc, \[\nabla u=q\sigma,\qquad
 q=\frac{2(x+1)}{R^2},\qquad
 q-1=\frac{1-r^2}{R^2}>0,\qquad p=(q-1)\sigma.\] In the elastic part \(p=0\) and \(\nabla u=\sigma\). Both fields have zero stress divergence. On \(r=1\), \(\phi=\theta/2\), so \(u\) and \(\sigma\) match and there is no interface jump. The boundary apex \((-1,0)\) has an integrable displacement gradient; \(u\in BV\) and \(\sigma\in H^1_{\mathrm{loc}}\).\par
Remove the closed disc with centre \(c=(0,21/20)\) and radius \(\rho=1/20\). It lies outside the plastic half-disc and is tangent to its boundary at \((0,1)\). The resulting \(\Omega\) is a bounded connected Lipschitz domain. Prescribe the actual trace of \(u\) on every boundary component. The displayed equilibrium and flow rule calibrate \(u\) as a minimizer, with no boundary plastic measure. Its \(\Omega_{\mathrm{pl}}\) is the upper unit half-disc. Its saturation set is the corresponding closed cap within \(\Omega\), with the single extreme point \((0,1)\) removed; this set is still convex. Thus (H) holds.\par
For \(0<\delta<3/5\), put \[P_\pm=(\pm\sqrt{\delta/2},\,1-\delta/2),\qquad
 M=(0,1-\delta/2).\] The points \(P_\pm\) lie in \(\Omega_{\mathrm{pl}}\), and the squared excess of their distance from \(c\) over \((\rho+\delta)^2\) is \[\delta\left(\frac9{20}-\frac{3\delta}{4}\right)>0.\] Their distances from the other boundary components also exceed \(\delta\), whereas the distance from \(M\) to the hole is \(\delta/2\). Hence \(P_\pm\in A_\delta\) but their midpoint \(M\notin A_\delta\). For example, \(\delta=1/50\) gives \(P_\pm=(\pm1/10,99/100)\) and the exact positive squared excess \(87/10000\). This example concerns Lemma 4.12’s hypotheses; it does not satisfy, or purport to refute, Theorem 5.1’s simply connected domain assumption.\par
Local repair of the null-sweep step. A compact noncharacteristic interface arc can be covered by balls \(B\Subset\Omega\) avoiding the exceptional points. Use \(A=\Omega_{\mathrm{pl}}\cap B\), which is convex. The stress is \(H^1(A)\) and extends continuously to \(\overline A\) after shrinking the ball; this uses Proposition 4.2 and interior regularity. On a shorter arc its characteristics meet the interface transversely. Merlet’s lemma therefore applies locally. For the appendix covering argument, choose disjoint small boundary intervals covering the null set; noncrossing characteristic strips are then disjoint. Arbitrary overlapping disc covers would not justify that disjointness.\par
The local sweep of an interface set of zero length has zero area. Every point farther along a swept characteristic lies at a rational flow time from an interior point in such a local sweep. On compact subsets of \(\Omega_{\mathrm{pl}}\), the field is locally Lipschitz, so the corresponding flow maps are locally Lipschitz. Countably many compact flow charts and rational times propagate the null-area conclusion. This supplies the step used to identify the continuous characteristic representative with the original \(BV\) function, without assuming convexity of \(\Omega\).\par
Repair of the Proposition 3.25 invocation. The needed statement can also be proved locally. Let \(K=\Omega_{\mathrm{pl}}\) and suppose a nondegenerate characteristic-boundary segment has a relative-interior point \(a\in\Omega\). Convexity of \(K\) gives an exact plastic half-ball at \(a\), with an elastic half-ball opposite it. Every cluster value of \(\sigma(x)\) as \(x\in K\) tends to \(a\) must be a unit normal to this flat boundary: otherwise a fixed short step along the limiting characteristic enters \(K\), and constancy of \(\sigma\) on chords would put \(a\) on an interior characteristic, contradicting its definition. Connectedness fixes the sign of the normal on a shorter subsegment.\par
The ambient \(H^1\) stress has matching traces on the two half-balls. Thus the elastic harmonic displacement has constant tangential trace and constant unit normal derivative there. Subtract the matching affine function. The remainder has zero Dirichlet and normal traces on the flat diameter, so its extension by zero is weakly harmonic in a full ball. Analyticity forces it to vanish. This would give \(|\nabla u|=1\) in the elastic half-ball, contradicting \(|\sigma|<1\). Openness excludes endpoints as well. Therefore no nondegenerate characteristic-boundary component meets \(\Omega\).\par
Only at most two isolated characteristic interface points remain. At such a point \(a\), characteristic chords approaching \(a\) have length tending to zero. Indeed, a nondegenerate limiting chord either contains an interior point, which would put \(a\) on its characteristic, or lies in a supporting boundary segment. In the second case choose a nearby noncharacteristic interface point on that segment. Proposition 4.2 and constancy on the approximating chords force its characteristic to be tangent to the supporting line, which cannot meet the interior of the convex set. Both alternatives are impossible.\par
The endpoints of the shrinking chords are noncharacteristic interface points and tend to \(a\). The locally repaired Lemma 4.12 identifies displacement along each chord with the continuous elastic trace at its endpoints. These values tend to the elastic trace at \(a\), furnishing the desired continuous extension. Continuous local representatives agree on overlapping open sweeps because they agree almost everywhere. This also covers the case where the fan-free component is all of \(K\); it does not assume an unavailable separating characteristic.\par
These arguments repair the cited continuity steps using the earlier regularity and characteristic results and the original hypotheses. They do not retroactively validate the printed derivation. Theorem 4.10 is a main new regularity result identified on source p. 6. The invalid convex truncation is needed to identify the continuous characteristic representative with the original \(BV\) solution; the unlicensed use of Proposition 3.25 is needed for the subsequent interface argument. These are essential dependencies of that original proof, so the finding is graded Major. The independently checked replacement restores the continuity conclusion; a false theorem or failure of the entire paper is not asserted or required for this classification.\par
\section*{Assessment and limits}
The report records one Major finding and one Minor finding. The examples include equilibrium, calibration, geometric and measure-normalization checks. The corrected arguments are provided as current reviewer derivations, not attributed to the authors or to historical checks.\par
Section 6’s no-return and no-closed-orbit arguments use Jordan interiors and should be read with their domain-topology scope explicit. On a simply connected domain, a global stream function \(\Psi\) with \(\nabla\Psi=\sigma^\perp\) gives \((\Psi\circ\gamma)'=|\sigma|^2\), ruling out nonconstant closed trajectories. The annular example in Section 3.3 has empty plastic interior and therefore is not a counterexample under the subsequent standing nonempty-interior assumptions. No contrary trajectory satisfying the full standing hypotheses was established, so this topology observation is not counted as a defect. Existence of the global continuous-field flow is also distinct from a general claim of uniqueness at exceptional points.\par
The external characteristic classification and imported regularity results remain explicit dependencies. The report verifies the identified proof structure and corrections in the acquired version; it is not an unconditional certification of every imported theorem. Human scientific adjudication remains pending.\par
9 J.-F. Babadjian and G. A. Francfort, Uniqueness, regularity, and characteristic flow for a non strictly convex singular variational problem, arXiv:2312.01849v2 (24 October 2024; manuscript dated 25 October 2024), 44 pp. Assigned journal DOI: https://doi.org/10.1002/cpa.70015.\par
\end{document}
